\section{CUDA 的长期赌注：技术资产如何变成战略资产}

\subsection{把 CUDA 放到 GeForce 的代价}

访谈里最有价值的历史片段之一是 Jensen 回顾 ``putting CUDA on GeForce''。他明确说这是当年 ``could not afford to do'' 的决定，短期财务压力很大，但长期看是必要动作。原因很直接：没有安装基数，就不会有开发者生态；没有生态，平台优势不可持续。

\begin{importantbox}{Installed Base 是第一性战略变量}
Jensen 直接把 CUDA 护城河归结为 installed base。不是某一代算子库领先，不是某个 benchmark 领先，而是 ``开发者心智 + 软件积累 + 迭代速度'' 的复合变量。
\end{importantbox}

\begin{figure}[H]
\centering
\includegraphics[width=0.92\textwidth]{frame-002.jpg}
\caption{访谈中段：回顾 CUDA 与 GeForce 的关键历史决策\protect\footnotemark}
\end{figure}
\footnotetext{视频画面时间区间：00:45:10--00:45:28。画面对应 Jensen 讨论生态扩展和系统路线。}

\subsection{为什么 OpenCL 没有形成同等飞轮}

Jensen 的观点不是否认替代技术的可行性，而是强调生态构建的时间函数。即使接口层面存在替代，开发者工具链、调优经验、框架适配、部署惯性会形成高迁移成本。这让 ``技术上可替代'' 和 ``商业上可替代'' 之间出现巨大差距。

\begin{warningbox}{平台竞争中的错觉}
很多团队把 ``语法兼容'' 当成 ``生态兼容''。实际上，生态迁移的核心成本来自性能调优、故障处理和工程经验沉淀，这些内容很少体现在 API 对照表里。
\end{warningbox}

\subsection{CUDA 进入 AI 时代后的新要求}

访谈后半段提到 CUDA 版本持续迭代（如 13.x 代际），并强调其价值不只在 training，也在 inference、agent pipeline、多模型协作等新场景。平台若想继续领先，需要在 ``通用性'' 与 ``特化能力'' 之间保持张力：既能承载新工作负载，又不能失去性能密度。

\begin{knowledgebox}{从 ``会用 CUDA'' 到 ``会运营 CUDA''}
在企业侧，CUDA 的真正门槛正在转向系统运营能力：如何做多租户隔离、推理弹性扩容、异构集群调度、故障退化和成本治理。开发能力与平台运营能力正在融合。
\end{knowledgebox}

\subsection{本章小结}

CUDA 的长期价值来自 ``先吃亏换规模'' 的战略路径。GeForce 承担了早期成本，installed base 形成后再反哺数据中心业务。这是一条典型的 ``生态先行，硬件兑现'' 路径。

